\documentclass[12pt]{article}
% Préambule commun à tous les documents extraits de « La Revue CAPES »
\usepackage[T1]{fontenc}
\usepackage[utf8]{inputenc}
\usepackage{lmodern}
\usepackage{amsmath,amssymb,amsthm,amsfonts,mathrsfs}
\usepackage{setspace}
\usepackage{listings}
\usepackage{pgfplots}
\pgfplotsset{compat=1.18}
\usepackage{longtable}
\usepackage{adjustbox}
\usepackage{lettrine}
\usepackage{tkz-tab}
\usepackage{changepage}
\usepackage{pdflscape}
\usepackage{graphicx}
\usepackage{tikz}
\usepackage{xcolor}
\usepackage[a4paper,top=2cm,bottom=2.5cm,left=2.5cm,right=2cm]{geometry}
\usepackage{fancyhdr}
\usepackage{draftwatermark}
\SetWatermarkText{La RevueCapes}
\SetWatermarkLightness{0.9}
\SetWatermarkScale{2}
\usepackage{hyperref}
\hypersetup{colorlinks=true,linkcolor=blue!50!black,urlcolor=blue!50!black}

\definecolor{revue}{RGB}{20,60,120}

\pagestyle{fancy}
\fancyhf{}
\renewcommand{\headrulewidth}{0.4pt}
\setlength{\headheight}{15pt}

% \entete{Matière}{Titre}{Type de document}
\newcommand{\entete}[3]{%
  \fancyhead[L]{\small\textsc{La Revue CAPES} -- #1}%
  \fancyhead[R]{\small #2}%
  \fancyfoot[C]{\thepage}%
  \thispagestyle{plain}%
  \begin{center}
    {\color{revue}\rule{\textwidth}{1.2pt}}\\[4pt]
    {\small\textsc{Concours directs CAP-CEG / CAPES -- Burkina Faso}}\\[6pt]
    {\LARGE\bfseries\color{revue} #2}\\[6pt]
    {\large\itshape #1 -- #3}\\[2pt]
    {\color{revue}\rule{\textwidth}{1.2pt}}
  \end{center}
  \vspace{0.5cm}%
}

% Encadré pour les remarques ajoutées dans les corrigés
\newcommand{\remarque}[1]{\par\medskip\noindent\fcolorbox{revue}{revue!6}{\parbox{\dimexpr\linewidth-2\fboxsep-2\fboxrule}{\textbf{\color{revue}Remarque.} #1}}\par\medskip}


\begin{document}
\entete{Mathématiques}{Corrigé détaillé -- Session 2020}{Correction}
\begin{spacing}{1.5}

\medskip\textbf{Exercice 1}

On donne la matrice $ A= \begin{pmatrix}
1 & 0 & -1\\
0 & 1 & 0\\
-1 & 1 & 1
\end{pmatrix}$                                                                   

1. Justifions que la matrice $A$ est inversible.

$det(A)=\begin{vmatrix}
1 & 0 & -1\\
0 & 1 & 0\\
-1 & 1 & 1
\end{vmatrix}=1-1=0$, alors $A$ n'est pas inversible.

2. Déterminons le polynôme caractéristique de $A$.

$P_A(x)=det(A-xI_3)=\begin{vmatrix}
1-x & 0 & -1\\
0 & 1-x & 0\\
-1 & 1 & 1-x
\end{vmatrix}=(1-x)\begin{vmatrix}
1-x & -1 \\ 
-1 & 1-x
\end{vmatrix}$

$=(1-x)((1-x)^2-1)=(1-x)(1-x-1)(1-x+1)=-x(1-x)(2-x)$ 

3. Soit $\alpha, \beta$ et $\gamma$ les valeurs propres de la matrice $A$. Déterminons $\alpha, \beta$ et $\gamma$ sachant qu'elles sont classées dans l'ordre croissant.

$\alpha=0$, $\beta=1$ et $\gamma=2$


4. Déterminons les vecteurs propres $u$, $v$ et $w$ associés aux valeurs propres $\alpha$, $\beta$ et $\gamma$.

Soit $E_\lambda$ le sous-espace propre lié à la valeur propre $\lambda$

\[
E_\lambda=\{u\in \mathbb{R}^3/Au=\lambda u\}
\]

\medskip
$\bullet$ Pour $\lambda=\alpha=0$, on a :

$\forall u\in E_\alpha$, \[Au=0\Rightarrow \begin{pmatrix}
1 & 0 & -1\\
0 & 1 & 0\\
-1 & 1 & 1
\end{pmatrix}\begin{pmatrix}
x\\
y\\
z
\end{pmatrix}=\begin{pmatrix}
0\\
0\\
0
\end{pmatrix}\Rightarrow \begin{cases}
x-z=0(1)\\
y=0(2)\\
-x+y+z=0(3)
\end{cases}\]

$(1)\Rightarrow x=z \hspace{0.7cm} (4)$

$(2)\Rightarrow y=0$

D'où $u=\begin{pmatrix}
x\\
y\\
z
\end{pmatrix}=\begin{pmatrix}
x\\
0\\
x
\end{pmatrix}=x\begin{pmatrix}
1\\
0\\
1
\end{pmatrix}\Rightarrow E_2=\langle \begin{pmatrix}
1\\
0\\
1
\end{pmatrix}\rangle$ et $u=\begin{pmatrix}
1\\
0\\
1
\end{pmatrix}$

\medskip
$\bullet$ Pour $\lambda=\beta=1$, on a :

$\forall v\in E_\beta$, \[Av=v\Rightarrow \begin{pmatrix}
1 & 0 & -1\\
0 & 1 & 0\\
-1 & 1 & 1
\end{pmatrix}\begin{pmatrix}
x\\
y\\
z
\end{pmatrix}=\begin{pmatrix}
x\\
y\\
z
\end{pmatrix}\Rightarrow \begin{cases}
x-z=x(1)\\
y=y(2)\\
-x+y+z=z(3)
\end{cases}\]

$(1)\Rightarrow z=0 \hspace{0.7cm} (4)$

$(3)\Rightarrow y=x$

D'où $u=\begin{pmatrix}
x\\
y\\
z
\end{pmatrix}=\begin{pmatrix}
x\\
x\\
0
\end{pmatrix}=x\begin{pmatrix}
1\\
1\\
0
\end{pmatrix}\Rightarrow E_\beta=\langle \begin{pmatrix}
1\\
1\\
0
\end{pmatrix}\rangle$ et $v=\begin{pmatrix}
1\\
1\\
0
\end{pmatrix}$

$\bullet$ Pour $\lambda=\gamma=2$, on a :

$\forall w\in E_\beta$, \[Aw=2w\Rightarrow \begin{pmatrix}
1 & 0 & -1\\
0 & 1 & 0\\
-1 & 1 & 1
\end{pmatrix}\begin{pmatrix}
x\\
y\\
z
\end{pmatrix}=2\begin{pmatrix}
x\\
y\\
z
\end{pmatrix}\Rightarrow \begin{cases}
x-z=2x\\
y=2y\\
-x+y+z=2z0
\end{cases}\]

$(1)\Rightarrow z=-x \hspace{0.7cm} (4)$

$(2)\Rightarrow y=0$

D'où $w=\begin{pmatrix}
x\\
y\\
z
\end{pmatrix}=\begin{pmatrix}
x\\
0\\
-x
\end{pmatrix}=x\begin{pmatrix}
1\\
0\\
-1
\end{pmatrix}\Rightarrow E_2=\langle \begin{pmatrix}
1\\
0\\
-1
\end{pmatrix}\rangle$ et $w=\begin{pmatrix}
1 \\
0 \\
-1
\end{pmatrix}$


5. Soit $D=\begin{pmatrix}
\alpha & 0 & 0 \\ 
0 & \beta & 0 \\ 
0 & 0 & \gamma
\end{pmatrix} $.

Donnons la matrice de passage $P$ de $A$ à $D$.

$P=\begin{pmatrix}
1 & 1 & 1  \\ 
0 & 1 & 0  \\ 
1 & 0 & -1
\end{pmatrix}$

\medskip
\textbf{Exercice 2}

1. On considère la fonction $f$ définie sur $\mathbb{R}$ par $f(x)=e^{-x}$.


a. Montrer que $f$ est dérivable sur $\mathbb{R}$.

$f$ est définie, continue et dérivable sur sont domaine de définition qui est $\mathbb{R}$


b. Déterminons $f^\prime(x)$, $f^{\prime\prime}(x)$ et $f^{(3)}(x)$; $\forall {x}\in \mathbb{R}$.

$\forall {x}\in \mathbb{R}$,

$f^\prime(x)=-e^{-x}$

$f^{\prime\prime}(x)=e^{-x}$

$f^{(3)}(x)=-e^{-x}$



c. Démontrons par récurrence que : $f^{(n)}(x)=(-1)^ne^{-x}$

-Pour $n=0$, $f^{(0)}(x)=(-1)^0e^{-x}=e^{-x}=f(x)$ (vraie)

-Pour $n=1$, $f^{(1)}(x)=(-1)^1e^{-x}=-e^{-x}=f'(x)$ (vraie)

-Pour $n=2$, $f^{(2)}(x)=(-1)^2e^{-x}=e^{-x}=f''(x)$ (vraie)

-Supposons que $f^{(n)}(x)=(-1)^ne^{-x}$ est vraie à l'ordre $n$ et vérifions à l'ordre $n+1$.

$f^{(n)}(x)=(-1)^ne^{-x}\Rightarrow (f^{(n)}(x))^\prime=((-1)^ne^{-x})^\prime=-(-1)^ne^{-x}=(-1)^{n+1}e^{-x}$ avec $(f^{(n)}(x))^\prime=f^{(n+1)}(x)$, alors $f^{(n+1)}(x)=(-1)^{n+1}e^{-x}$, d'où la propriété est vraie à l'ordre $n+1$.

On conclut donc que $f^{(n)}(x)=(-1)^ne^{-x}$, $\forall n\in \mathbb{N}$

d. Donnons le développement limité d'ordre $n$ de $f$ au voisinage de $0$.

De la formule de Taylor, le développement limité d'ordre $n$ d'une fonction $f$ au voisinage de $x_0$ est donné par :

\[
f(x)=\sum_{k=0}^{n}\frac{(x-x_0)^k}{k!}f^{(k)}(x_0)+\tau(x^n)
\]

Dans notre cas, $x_0=0$

$f^{(n)}(x)=(-1)^ne^{-x}\Rightarrow f^{(n)}(0)=(-1)^ne^{-0}=(-1)^n$

Donc,
\begin{align*}
f(x)&=\sum_{k=0}^{n}\frac{(x-0)^k}{k!}f^{(k)}(0)+\circ(x^n)\\
&=\sum_{k=0}^{n}\frac{x^k}{k!}(-1)^k+\circ(x^n)\\
&=\sum_{k=0}^{n}(-1)^k\frac{x^k}{k!}+\circ(x^n)\\
f(x)&=1-\frac{x}{1!}+\frac{x^2}{2!}+\cdots (-1)^n\frac{x^n}{n!}+\circ(x^n)
\end{align*}

2. On considère la courbe paramétrée : $(\Gamma) :\begin{cases}
x(t)=\cos 3t\\
y(t)=\sin 2t
\end{cases} \forall t\in \mathbb{R}$

a. Montrons que l'étude de la courbe peut être réduite à l'intervalle $[0,\pi]$.

\textbf{Vérifions la périodicité}

On sait que $\cos(at)$ est périodique de période $\frac{2\pi}{\vert a\vert}$, et $\sin(bt)$ est périodique de période $\frac{2\pi}{\vert b\vert}$

Donc, $\cos(3t)$ est périodique de période $\frac{2\pi}{3}$ et $\sin(2t)$ est périodique de période $\pi$

Le $PPCM(\frac{2\pi}{3},\pi)=2\pi$ , donc la courbe est périodique de période $2\pi$. Ainsi, le domaine d'étude de la courbe peut être réduit à un intervalle d'amplitude $2\pi$, soit $[-\pi,\pi]$

\medskip
\textbf{Vérifions la parité}

$\begin{cases}
x(-t)=\cos(-3t)=\cos 3t=x(t)\\
y(-t)=\sin(-2t)=-\sin 2t=-y(t)
\end{cases}$

$x(t)$ est paire et $y(t)$ est impaire, alors la courbe est symétrique par rapport à l'axe des abscisses. Donc le domaine d'étude de la courbe peut être réduit à $[0,\pi]$.

b. Dressons le tableau de variation commun de $x$ et $y$ sur $[0,\pi]$.

Étudions la courbe.

\textbf{Étude de $x(t)$}

$x(t)=\cos 3t$ est définie, continue et dérivable sur $\mathbb{R}$, car fonction trigonométrique, et

$x'(t)=-3\sin 3t$

\medskip
Signe de $x'(t)$

soit $x'(t)=-3\sin 3t=0\Rightarrow \sin 3t=0=\sin 0\Rightarrow \begin{cases}
3t=2k\pi\\
3t=\pi+2k\pi
\end{cases}\Rightarrow \begin{cases}
t=\frac{2k\pi}{3}\\
t=\frac{\pi}{3}+\frac{2k\pi}{3}
\end{cases}\forall k\in \mathbb{Z}$

-Pour $k=0, t=0$ ou $t=\frac{\pi}{3}$

-Pour $k=1,t=\frac{2\pi}{3}$ ou $t=\pi$

\begin{tikzpicture}
  \tkzTabInit[lgt=3, espcl=3]{$t$ /1, $x'(t)$ /1}{$0$, $\frac{\pi}{3}$, $\frac{2\pi}{3}$, $\pi$}
  \tkzTabLine{, -, z, +, z, -,}
\end{tikzpicture}


\underline{Variations de $x(t)$}

\[
\begin{cases}
\forall t\in [0,\frac{\pi}{3}],x'(t)\leq 0\hspace*{0.5cm} alors\hspace*{0.5cm} x(t)\hspace*{0.5cm} est \hspace*{0.5cm}decroissante\\
\forall t\in [\frac{\pi}{3}, \frac{2\pi}{3}], x'(t)\geq 0 \hspace*{0.5cm} alors\hspace*{0.5cm} x(t)\hspace*{0.5cm} est\hspace*{0.5cm} croissante\\
\forall t\in [\frac{2\pi}{3},\pi], x'(t)\leq 0\hspace*{0.5cm} alors\hspace*{0.5cm} x(t)\hspace*{0.5cm} est\hspace*{0.5cm} decroissante
\end{cases}
\]

\medskip\textbf{Étude de $y(t)$ sur $[0,\pi]$}

$y(t)=\sin{2t}\Rightarrow y'(t)=2\cos 2t$

$y'(t)=0\Rightarrow \cos 2t=0=\cos\frac{\pi}{2}$

\[
\begin{cases}
2t=\frac{\pi}{2}+2k\pi\\
2t=-\frac{\pi}{2}+2k\pi
\end{cases}\Rightarrow
\begin{cases}
t=\frac{\pi}{4}+k\pi\\
t=-\frac{\pi}{4}+k\pi
\end{cases}
\]

-Pour $k=0$, $t=\frac{\pi}{4}$

-Pour $k=1$, $t=\frac{3\pi}{4}$

Tableau de signe de $y'(t)$

\begin{tikzpicture}
  \tkzTabInit[lgt=3, espcl=3]{$t$ /1, $y'(t)$ /1}{$0$, $\frac{\pi}{4}$, $\frac{3\pi}{4}$, $\pi$}
  \tkzTabLine{, +, z, -, z, +,}
\end{tikzpicture}

\underline{Variations de $y(t)$}

\[
\begin{cases}
\forall t\in [0,\frac{\pi}{4}],y'(t)\geq 0\hspace*{0.5cm} alors\hspace*{0.5cm} y(t)\hspace*{0.5cm} est \hspace*{0.5cm}croissante\\
\forall t\in [\frac{\pi}{4}, \frac{3\pi}{4}], y'(t)\leq 0 \hspace*{0.5cm} alors\hspace*{0.5cm} y(t)\hspace*{0.5cm} est\hspace*{0.5cm} decroissante\\
\forall t\in [\frac{3\pi}{4},\pi], y'(t)\leq 0\hspace*{0.5cm} alors\hspace*{0.5cm} y(t)\hspace*{0.5cm} est\hspace*{0.5cm} croissante
\end{cases}
\]

\medskip\textbf{Tableau commun de variation des fonctions $x$ et $y$}

$x(0)=1$, $x(\frac{\pi}{4})=-\frac{\sqrt{2}}{2}$, $x(\frac{\pi}{3})=-1$, $x(\frac{2\pi}{3})=1$, $x(\frac{3\pi}{4})=-1$, $x(\pi)=-1$

$x'(0)=0$, $x'(\pi)=0$

$y(0)=0$, $y(\frac{\pi}{4})=1$, $y(\frac{\pi}{3})=\frac{\sqrt{3}}{2}$, $y(\frac{2\pi}{3})=-\frac{\sqrt{3}}{2}$, $y(\frac{3\pi}{4})=-1$, $y(\pi)=0$

$y'(0)=2$, $y'(\pi)=2$

\begin{center}
\resizebox{\textwidth}{!}{\begin{tikzpicture}
    % Initialisation du tableau
    \tkzTabInit[lgt=2, espcl=3]
    {$t$ /1, $x'(t)$ /1, $x(t)$ /3, $y'(t)$ /1, $y(t)$ /3} 
    {$0$, $\frac{\pi}{4}$, $\frac{\pi}{3}$, $\frac{2\pi}{3}$, $\frac{3\pi}{4}$, $\pi$}

    % Ligne de signes de x'(t)
    \tkzTabLine{, ,-, ,z, +,z, ,- ,}
    
    
    % Ligne de variations de x(t)
    \tkzTabVar{+/$1$, R/ , -/$-1$, +/$1$, R/, -/$-1$,}
    
    %Pour mettre des valeurs(les images) sur les flèches
    \tkzTabIma{1}{3}{2}{$-\frac{\sqrt{2}}{2}$}
    
    \tkzTabIma{4}{6}{5}{$\frac{\sqrt{2}}{2}$}
    % Ligne de signes de y'(t)
    
    % Ligne de signes de y'(t)
    \tkzTabLine{,+,z, , ,- , , ,z,+}
    
    %Pour ajouter les valeurs de dérivées(vecteurs directeurs)
    
    
    % Ligne de variations de y(t)
    \tkzTabVar{-/$0$, +/$1$, R/ , R/ , -/$-1$, +/$0$,}
    
    %Pour mettre des valeurs (les images) sur les flèches.
    \tkzTabIma{2}{5}{3}{$\frac{\sqrt{3}}{2}$}
    
    \tkzTabIma{2}{5}{4}{$-\frac{\sqrt{3}}{2}$}
\end{tikzpicture}
}
\end{center} 

\medskip
\textbf{Exercice 3}

Soient $n\geq 1$, un entier naturel et $\mathbb{R}_n[X]$ l'espace vectoriel des polynômes réels de degré inférieur ou égal à $n$. $\forall {P,Q\in \mathbb{R}_n[x]}$, on pose :

\[ \beta(P,Q)=\int_{0}^{1}xP(x)Q^{\prime}(x). dx\]

1. Montrons que $\beta$ est une forme bilinéaire.

\underline{Montrons que $\beta$ est linéaire par rapport à la première variable}

Soient $P,Q$ et $G$ $\in \mathbb{R}_n[x]$ et $\alpha\in \mathbb{R}$

On a :

\begin{align*}
\beta(P+\alpha G,Q) &=\int_{0}^{1}x(P(x)+\alpha G(x)) Q^{\prime}(x).dx\\
&=\int_{0}^{1}(xP(x)Q^{\prime}(x)+\alpha xG(x)Q^{\prime}(x)).dx\\
&=\int_{0}^{1}xP(x)Q^{\prime}(x).dx+\alpha\int_{0}^{1}xG(x)Q^{\prime}(x).dx\\
\beta(P+\alpha G,Q) &=\beta(P,Q)+\alpha\beta(G,Q)
\end{align*}, alors $\beta$ est linéaire par rapport à la première variable.

 \medskip
 \underline{Montrons que $\beta$ est linéaire par rapport à la deuxième variable}
 
 Soient $P,Q$ et $G$ $\in \mathbb{R}_n[x]$ et $\alpha\in \mathbb{R}$
 
 On a :

\begin{align*}
\beta(P,\alpha G+Q)&=\int_{0}^{1}xP(x)(\alpha G + Q)^{\prime}(x). dx\\
&=\int_{0}^{1}xP(x)(\alpha G^{\prime}(x) + Q^{\prime}(x)). dx\\
&=\alpha\int_{0}^{1}xP(x)G^{\prime}(x).dx+\int_{0}^{1}xP(x)Q^{\prime}(x). dx\\
\beta(P,\alpha G+Q)&=\alpha\beta(P,G)+\beta(P,Q)
\end{align*}, alors $\beta$ est linéaire par rapport à la deuxième variable.

On conclut donc que $\beta$ est \textbf{bilinéaire}.
 

2. Exprimons $\beta(x,1)$ et $\beta(1,x)$.

$\beta(x,1)=\int_{0}^{1}x\times x\times 1^{\prime}. dx=\int_{0}^{1}x^2\times 0.dx =0$

$\beta(1,x)=\int_{0}^{1}x\times 1\times(x)^{\prime}. dx=\int_{0}^{1}x=[\frac{1}{2}x^2]_0^1=\frac{1}{2}$

3. $\beta$ est-elle symétrique? Antisymétrique?

$\beta$ est symétrique si $\beta(P,Q)=\beta(Q,P)$.

On a : $\beta(x,1)\neq\beta(1,x)$ alors $\beta$ n'est pas symétrique, mais antisymétrique.


\medskip
\textbf{Exercice 4}

1. Montrer à l'aide du théorème de Fubini que : 

$I=\int_{0}^{+\infty}e^{-\frac{x^2}{2}}dx=\sqrt{\frac{\pi}{2}}$

$I=\int_{0}^{+\infty}e^{-\frac{x^2}{2}}dx\Longrightarrow I^2=(\int_{0}^{+\infty}e^{-\frac{x^2}{2}}dx)^2=(\int_{0}^{+\infty}e^{-\frac{x^2}{2}}dx)\times (\int_{0}^{+\infty}e^{-\frac{y^2}{2}}dy)=\int_{0}^{+\infty}\int_{0}^{+\infty}e^{-\frac{x^2}{2}}\times e^{-\frac{y^2}{2}}dxdy=\int_{0}^{+\infty}\int_{0}^{+\infty}e^{-\frac{x^2+y^2}{2}}dxdy$

\medskip
Posons :

$\begin{cases}
x=r\cos\theta\\
y=r\sin\theta
\end{cases}$

$x^2+y^2=r^2$

\medskip
Calculons le Jacobien

$J(r,\theta)=\begin{vmatrix}
\frac{\partial x}{\partial r}&\frac{\partial x}{\partial \theta}\\
\frac{\partial y}{\partial r} & \frac{\partial y}{\partial \theta}
\end{vmatrix}=\begin{vmatrix}
\cos\theta & -r\sin\theta\\
\sin\theta & r\cos\theta
\end{vmatrix}=r$ $\Rightarrow dxdy=rdrd\theta$

Si $x\longrightarrow 0$ et $y\longrightarrow 0$ alors $r\longrightarrow 0$

Si $x\longrightarrow +\infty$ et $y\longrightarrow +\infty$ alors $r\longrightarrow +\infty$, mais de sorte que $cos\theta$ et $sin\theta$ soient positifs( donc $\theta\in [0,\frac{\pi}{2}]$).

Ainsi $r\in [0, +\infty[$ et $\theta\in [0,\frac{\pi}{2}]$

Donc, 

\[
I^2=\int_0^{+\infty}\int_0^{\frac{\pi}{2}}e^{-\frac{r^2}{2}}rdrd\theta=\int_0^{+\infty}re^{-\frac{r^2}{2}}dr\times \int_0^{\frac{\pi}{2}}d\theta=[-e^{-\frac{r^2}{2}}]_0^{+\infty}\times [\theta]_0^{\frac{\pi}{2}}=1\times \frac{\pi}{2}=\frac{\pi}{2}\]

$I^2=\frac{\pi}{2}$ Comme $I\geq 0$ alors $I=\sqrt{\frac{\pi}{2}}$


2. On pose $\forall {n\in \mathbb{N}}$, $\forall {x\in \mathbb{R}}$ $\\
f_n(x)=(1+\frac{x^2}{n})-\frac{n+1}{2}$ et $a_n=\int_{-\infty}^{+\infty}f_n(x)dx.$

a. Montrons que, $\forall {u\in \mathbb{R}}$ et $\forall {p\in \mathbb{N}}, (1+u)^p\geq 1+pu$.

Procédons par récurrence :

-Pour $p=1$, on a $(1+u)^1\geqslant 1+u$ (vrai)

-Pour $p=2$, on a $(1+u)^2=1+2u+u^2\geq 1+2u$ (vrai)

-Supposons que $(1+u)^p\geq 1+pu$ est vrai à l'ordre $p$, et vérifions à l'ordre $p+1$.

On a :

$(1+u)^p\geq 1+pu\Rightarrow (1+u)^p(1+u)\geq (1+pu)(1+u)\Rightarrow (1+u)^{p+1}\geq 1+pu + u+pu^2\Rightarrow (1+u)^{p+1}\geq 1+pu + u+pu^2$ avec $1+pu + u+pu^2=1+(1+p)u+pu^2\geq 1+(1+p)u$.

$\Rightarrow (1+u)^{p+1}\geq 1+(1+p)u$, d'où la propriété est vraie à l'ordre $p+1$.

On conclut donc que la propriété est vraie $\forall {u\in \mathbb{R}}$ et $\forall {p\in \mathbb{N}}$.


b. Déduisons que $\forall {x\in \mathbb{R}}, f_n(x) \leq \frac{2}{1+x^2}$.

On a  $f_n(x)=(1+\frac{x^2}{n})-\frac{n+1}{2}$ 

$(1+u)^p\geq 1+pu\Longrightarrow (1+u)^{-p}\leq \frac{1}{1+pu}$

De plus, $\frac{1}{1+pu}\leq \frac{2}{1+pu}\Longrightarrow (1+u)^{-p}\leq \frac{2}{1+pu}$

Posons $p=\frac{n+1}{2}$ et $u=\frac{x^2}{n}$

En reportant dans l'équation précédente, on obtient : 

$f_n(u)=(1+u)^{-p}$ et $f_n(x)\leq\frac{2}{1+x^2}$

c. Calculons $lim_{n\to +\infty}a_n$.

$lim_{n\to +\infty}a_n=lim_{n\to +\infty}\int_{-\infty}^{+\infty}f_n(x)dx\leq \lim_{n\to +\infty}\int_{-\infty}^{+\infty}\frac{2}{1+x^2}=\lim_{n\to +\infty}[2\arctan(x)]_{-\infty}^{+\infty}=\lim_{n\to +\infty}2(\frac{\pi}{2}+\frac{\pi}{2})=2\pi$

$\Longrightarrow lim_{n\to +\infty}a_n=2\pi$

\medskip

\end{spacing}
\end{document}
